\chapter{TRABALHOS RELACIONADOS}
\label{cap:trabalhos_relacionados}

A crescente informatização das organizações rurais tem impulsionado o desenvolvimento de sistemas voltados à gestão de associações, cooperativas e atividades produtivas. No contexto da apicultura, diferentes soluções vêm sendo propostas para apoiar tanto a administração das organizações quanto o gerenciamento de informações relacionadas à produção. Embora esses trabalhos contribuam para a modernização do setor, eles apresentam enfoques distintos, priorizando aspectos específicos da gestão ou da atividade apícola. 


\section{APISAPP}
\label{sec:apisapp}

Pinto (2016) desenvolveu o APISAPP, uma solução que integra um aplicativo móvel para a coleta de dados em campo a um website de suporte à gestão do apiário. O sistema foi concebido com foco na viabilidade econômica da atividade apícola, tratando o apiário como um empreendimento rural e oferecendo ferramentas de planejamento financeiro, como fluxo de caixa, cálculo de custos de produção e definição do preço de venda do mel. Além disso, a solução apresenta caráter educativo ao disponibilizar planilhas de produção projetada, como o calendário alimentar e projeções de produtividade, permitindo ao apicultor compreender o impacto de práticas de manejo, como alimentação artificial e substituição de rainhas, sobre o desempenho das colmeias. Para sua implementação, foram empregadas ferramentas de baixo código, como Wix e Fábrica de Aplicativos, além do uso de formulários por meio do SurveyMonkey.

Embora o APISAPP contribua para a gestão e para a compreensão da atividade apícola sob as perspectivas produtiva e econômica, sua abordagem permanece centrada no apicultor individual e na unidade produtiva isolada, sem contemplar os processos de coordenação e governança coletiva característicos das associações. Assim, trata-se de uma ferramenta voltada ao apoio da tomada de decisão no nível do produtor, ainda que inserido em um contexto associativo, sem atuar diretamente sobre a estrutura organizacional da entidade coletiva.

O presente trabalho, por sua vez, direciona-se justamente a esse nível organizacional ao propor a automatização dos processos internos de uma associação de apicultores. A solução busca enfrentar problemas relacionados à informalidade, à baixa padronização dos processos e à ausência de mecanismos estruturados de transparência, fatores que comprometem a gestão coletiva. Dessa forma, enquanto o APISAPP procura apoiar a tomada de decisão no âmbito individual da produção, a proposta apresentada fundamenta-se nos princípios da Economia Solidária.

Além das diferenças quanto ao domínio de aplicação, observa-se também uma distinção no processo de desenvolvimento adotado. Enquanto Pinto (2016) prioriza a viabilidade da solução por meio de ferramentas de baixo código, este trabalho fundamenta-se em práticas da Engenharia de Software e utilizando tecnologias como Go e Next.js.

\section{Colonymon Web}
\label{sec:colonymon}

Maciel (2024) apresenta o desenvolvimento do Colonymon Web, uma plataforma voltada para a gestão técnica de apiários com ênfase na padronização de inspeções biológicas. O trabalho destaca-se pela adoção de um processo estruturado de Engenharia de Requisitos, conduzido por meio de reuniões com especialistas com mais de 30 anos de experiência na atividade apícola para a elaboração de um questionário de inspeção detalhado. A aplicação permite registrar informações relacionadas ao estado das colmeias, como força da colônia, presença de mel verde e condições da rainha, fornecendo uma visão técnica do apiário que auxilia o acompanhamento da produção.

Embora o Colonymon Web tenha como foco o gerenciamento técnico dos apiários, o sistema incorpora uma dimensão coletiva ao permitir que diferentes usuários compartilhem um mesmo apiário e dividam as responsabilidades relacionadas ao manejo das colmeias. Essa funcionalidade reconhece que a atividade apícola pode ser desenvolvida de forma colaborativa, superando a perspectiva do produtor que atua de maneira isolada. Entretanto, essa colaboração permanece restrita ao nível operacional da produção, uma vez que o sistema não se destina à gestão de uma organização coletiva. Em contraste, a proposta deste trabalho direciona-se à administração de uma associação de apicultores, na qual a colaboração ultrapassa o compartilhamento das atividades de campo e passa a envolver processos de autogestão, transparência e coordenação organizacional, ou seja, fundamentados nos princípios da Economia Solidária.

Também se observa uma diferença no processo de desenvolvimento adotado. Maciel (2024) emprega práticas de Engenharia de Software, incluindo a modelagem por meio de diagramas UML e a utilização de tecnologias modernas como React e Next.js. Entretanto, a avaliação do sistema baseia-se predominantemente em testes manuais. Neste trabalho, por sua vez, adota-se o Test-Driven Development (TDD) como estratégia de Verificação e Validação, buscando garantir que as regras de negócio permaneçam corretas ao longo da evolução do sistema.


\section{SisFarming}
\label{sec:sisfarming}

Aguiar (2026) propõe o SisFarming, um ecossistema digital composto por um sistema desktop voltado à gestão centralizada de associações da agricultura familiar e por um protótipo mobile destinado ao uso direto pelos produtores. Entretanto, enquanto a abordagem de Aguiar (2026) foca na Teoria da Agência sob a ótica clássica da mitigação de conflitos entre Agente e Principal, a solução proposta neste presente trabalho submete essa relação aos princípios da Economia Solidária. Nesse sentido, a ferramenta proposta aqui não é vista apenas como um instrumento de controle de interesses divergentes, mas como uma garantia de que a cooperação e a ajuda mútua sejam preservadas, transformando a transparência em um pilar de sustentação para a autogestão e para a coesão do grupo.

Apesar dessa diferença de fundamentação teórica, o SisFarming aproxima-se da proposta deste trabalho por tratar a associação como uma instituição que necessita de mecanismos de transparência para sustentar a confiança entre seus membros. Contudo, Aguiar (2026) relata que o protótipo mobile do sistema não alcançou adoção operacional efetiva entre os produtores, atribuindo essa dificuldade a restrições cognitivas e organizacionais dos usuários frente à ferramenta proposta. Esse resultado reforça a importância de uma elicitação de requisitos conduzida de forma próxima ao público-alvo, o que, no presente trabalho, é endereçado por meio da Prototipação Exploratória, técnica escolhida justamente por se adequar a stakeholders com baixa familiaridade com ferramentas digitais. Além disso, o trabalho de Aguiar (2026) não menciona a adoção do Test-Driven Development ou de qualquer outra prática de verificação automatizada como parte do processo de desenvolvimento, o que distancia sua abordagem da preocupação com a confiabilidade do software.


\section{Gestão de Associações}
\label{sec:gestao_associacoes}

Oliveira et al. (2021) apresentam o desenvolvimento de um sistema voltado à informatização de atividades administrativas em associações de caráter genérico. O trabalho destaca-se por ser uma aplicação prática da Engenharia de Software, fundamentada no uso de metodologias ágeis como Scrum e Extreme Programming (XP). Os autores enfatizam o uso da técnica de programação em pares, uma prática central do XP, como estratégia para promover a colaboração entre os desenvolvedores e contribuir para a qualidade do software em ambientes de desenvolvimento acadêmico. Em termos de arquitetura, o sistema utiliza uma API independente, persistência de dados em MongoDB e comunicação via JSON, garantindo a independência entre os módulos e facilitando manutenções futuras.

A proposta de Oliveira et al. (2021) converge com o presente trabalho ao reconhecer que a ausência de práticas sistemáticas de garantia da qualidade favorece o desenvolvimento de produtos que podem não atender plenamente aos requisitos dos usuários. Contudo, enquanto o trabalho anterior concentra seus esforços na colaboração interpessoal por meio da programação em pares para promover essa qualidade, a proposta apresentada adota o Test-Driven Development (TDD) como pilar metodológico central. Embora o XP inclua diretrizes de teste, este trabalho aprofunda as práticas de Verificação e Validação (V\&V) ao implementar o ciclo Red-Green-Refactor, assegurando que cada regra de negócio da Associação de Apicultores de Poranga seja formalmente validada por testes automatizados antes de sua implementação definitiva.

Além do diferencial metodológico proporcionado pela adoção do TDD, a solução proposta neste trabalho também se distingue da apresentada por Oliveira et al. (2021) ao adotar a estratégia Mobile First. Enquanto os autores direcionam sua solução para portais administrativos de uso predominantemente web, a aplicação desenvolvida utiliza tecnologias como Go e Next.js para entregar uma solução simultaneamente robusta do ponto de vista lógico e acessível à realidade dos produtores rurais.

\section{Análise Comparativa}
\label{sec:analise_comparativa}

A análise dos trabalhos relacionados evidencia que as soluções existentes atendem a necessidades específicas da apicultura ou da gestão de organizações, porém sob perspectivas distintas. Os trabalhos de Pinto (2016) e Maciel (2024) concentram-se no apoio às atividades produtivas e ao manejo apícola, oferecendo funcionalidades voltadas ao planejamento da produção e ao acompanhamento técnico dos apiários. Por outro lado, Oliveira et al. (2021) e Aguiar (2026) direcionam seus esforços para a gestão institucional de associações, priorizando processos administrativos e organizacionais sem contemplar particularidades do contexto apícola.

Sob a perspectiva da Engenharia de Software, observa-se ainda que os trabalhos analisados empregam diferentes metodologias e tecnologias de desenvolvimento, porém com pouca ênfase em estratégias que garantam maior confiabilidade ao software. Em especial, Maciel (2024) restringe a validação do sistema à realização de testes manuais, enquanto os demais trabalhos não apresentam uma abordagem baseada em testes automatizados orientados ao desenvolvimento. Nesse contexto, este trabalho diferencia-se ao adotar o Test-Driven Development (TDD), incorporando a verificação automatizada desde as etapas iniciais da implementação e contribuindo para a confiabilidade das regras de negócio do sistema.

Outro aspecto que diferencia a proposta desenvolvida refere-se à integração entre a gestão institucional e as necessidades específicas da apicultura. Enquanto os sistemas voltados à produção concentram-se no gerenciamento individual ou técnico dos apiários e os sistemas para associações apresentam soluções de caráter mais genérico, este trabalho busca reunir essas duas perspectivas em uma única aplicação, considerando simultaneamente as demandas administrativas de uma associação apícola e as características do seu contexto produtivo. Além disso, a adoção da estratégia Mobile First busca ampliar a acessibilidade da aplicação em dispositivos móveis, realidade predominante entre os produtores rurais, reforçando o alinhamento da solução às condições de uso dos associados.

\begin{table}[h]
\centering
\caption{Comparação entre os trabalhos relacionados}
\label{tab:comparativo}
\begin{tabular}{|l|c|c|c|c|c|}
\hline
\textbf{Critério} & \textbf{SisFarming} & \textbf{APISAPP} & \textbf{Colonymon} & \textbf{Oliveira et al.} \\
& \textbf{(2026)} & \textbf{(2016)} & \textbf{(2024)} & \textbf{(2021)} \\
\hline
Domínio específico & Agric. Familiar & Apicultura & Apicultura & Associações \\
\hline
Foco principal & Institucional & Individual & Produtiva & Institucional \\
\hline
Metodologia de V\&V & Não especificado & Manual & Manual & Verificação ágil \\
\hline
Estratégia de interface & Híbrida & Mobile/Web & Web & Web \\
\hline
Base tecnológica & Pascal/Delphi & Low-code & React/Next.js & API/MongoDB \\
\hline
Economia Solidária & Não & Não & Não & Não \\
\hline
\end{tabular}
\end{table}